% =====================================================================
% Methodology, Results, Discussion, Limitations, Conclusion, Future Work
%
% Placeholders use \NUM{} for single figures and \TODO{} for prose that
% depends on the final n=250 run. Add this to your preamble:
%
%   \usepackage{graphicx}
%   \graphicspath{{figures/}{figs/}{figured/}{./}}
%   \usepackage{xcolor}
%   \newcommand{\NUM}[1]{\textcolor{red}{\textbf{[#1]}}}
%   \newcommand{\TODO}[1]{\textcolor{red}{[TODO: #1]}}
%
% Everything NOT wrapped in \NUM or \TODO is final and does not depend
% on the running evaluation.
% =====================================================================

\section{Methodology}

\begin{figure*}[t]
\centering
\includegraphics[width=\textwidth]{fig_method.png}
\caption{The study in full. Decontamination against three public Bengali
mathematics benchmarks precedes difficulty filtering, so the LIMO window is
applied to an already-clean pool. The two training sets differ only in the
difficulty band they are drawn from. Every backbone receives the identical
LoRA configuration and is evaluated on the identical held-out items in both
conditions, which makes every comparison in this paper a paired one.}
\label{fig:method}
\end{figure*}

\subsection{Experimental Design}

The study is a controlled experiment with two independent variables, parameter
scale and targeted Bengali pretraining depth, and one intervention, LIMO-style
minimal-data fine-tuning. Every model is evaluated twice on the same held-out
problems, once as released and once after fine-tuning, so each model acts as its
own control.

Seven backbones were selected so that the two independent variables move
separately. Table~\ref{tab:models} summarises them. TigerLLM-1B and TituLLM-1B
are close in size but differ by roughly $3{,}700\times$ in targeted Bengali
pretraining. TituLLM-1B and TituLLM-3B share a pretraining recipe and differ
only in size. The Qwen3 family spans 0.6B to 4B with no Bengali-specific
pretraining, isolating scale.

Two of the Bengali models are available in both base and instruction-tuned
form. Because TigerLLM is instruction-tuned and TituLLM's base models are not, a
comparison between them would confound pretraining depth with instruction
tuning. We therefore evaluate the instruction-tuned TituLLM variants as the
primary depth comparison and retain the base variants as an instruction-tuning
ablation.

\begin{table}[htbp]
\caption{The seven backbones under test. Bengali depth records whether the model received Bengali-specific continued pretraining beyond a general multilingual mix. Parameter counts are total, not trainable; LoRA trains roughly 0.5\% of them.}
\label{tab:models}
\centering
\begin{tabular}{|l|l|c|c|}
\hline
\textbf{Model} & \textbf{Base} & \textbf{Bengali PT} & \textbf{Params} \\
\hline
TigerLLM-1B-it        & Gemma3      & $\sim$10M tokens & 1.0B \\
TituLLM-1B v2.0       & Llama-3.2   & $\sim$37B tokens & 1.2B \\
TituLLM-1B-Instruct   & Llama-3.2   & $\sim$37B tokens & 1.2B \\
TituLLM-3B v2.0       & Llama-3.2   & $\sim$37B tokens & 3.2B \\
TituLLM-3B-Instruct   & Llama-3.2   & $\sim$37B tokens & 3.2B \\
Qwen3-0.6B            & Qwen3       & incidental       & 0.6B \\
Qwen3-1.7B            & Qwen3       & incidental       & 1.7B \\
Qwen3-4B              & Qwen3       & incidental       & 4.0B \\
\hline
\end{tabular}
\end{table}

\subsection{Training Data Curation}

Training data was drawn from the SFT split of the Ganit dataset
\parencite{ganitllm}, a Bengali mathematical reasoning corpus of 16,868
problems, each carrying a full Bengali reasoning trace and a difficulty score.
The difficulty score is the number of times Qwen3-32B solved the problem across
32 sampled attempts, so it records difficulty relative to a 32B model. The pool
is assembled from four upstream sources: numina-math-cot-bn (13,160 problems),
SOMADHAN (3,487), mCoT-MATH-bn (193) and s1k-32-Bangla (28). Its own difficulty
tags skew easy (10,065 easy, 6,411 olympiad, 258 hard, 134 medium), but the
solve-rate column tells a harsher story: 5,645 problems, a third of the corpus,
were never solved by Qwen3-32B in any of 32 attempts, and 91.7\% of that
unsolvable band comes from numina-math-cot-bn alone. We treat the zero band as
unreliable rather than merely hard, and exclude it.

Figure~\ref{fig:method} shows the pipeline. Decontamination ran first, before
any difficulty filtering, against 4,080 problems pooled from three public
Bengali mathematics benchmarks: Bn-MGSM (250 problems), Bn-MSVAMP (1,000) and
BenNumEval (2,830 across five task types). A training problem was removed if it
matched a benchmark problem on either of two signals. The first is a numeric
fingerprint, the hash of the sorted set of numbers appearing in the normalised
problem text. Numbers survive re-translation even when every surrounding word
changes, so this catches exactly the case that n-gram or MinHash deduplication
misses. The second is token Jaccard overlap at a threshold of 0.7, which
catches near-identical wording. Normalisation maps Bengali numerals to Western
digits, lowercases, strips punctuation and collapses whitespace, so the two
numeral conventions present in the corpus compare on equal footing. Together
these removed 3,315 problems and left 13,553.

Following LIMO's selection criterion, we then retained problems in the window
$1 \leq \text{correct\_counts} \leq 16$, which reduced 13,553 to 873. A final
near-duplicate pass within the surviving set removed nothing, which is itself
mild evidence that the pool is not padded with reworded repeats. The resulting
set is heavily weighted towards hard problems: 74.6\% carry an olympiad
difficulty tag and 25.3\% are tagged hard, with a mean pass rate of 3.3 out of
32.

A second training set was constructed later to test whether that difficulty
distribution explains our results. Holding every other choice constant, we
selected 1,000 traces from the band $\text{correct\_counts} \geq 28$, giving a
mean pass rate of 31.8 out of 32 and a difficulty composition of 100\% easy.
Section~\ref{sec:difficulty} reports that experiment. This set was
decontaminated directly against the held-out evaluation items rather than
against the public benchmarks, using the same numeric fingerprint idea. That
filter is deliberately aggressive and removed 1,782 candidates.

The evaluation set was independently checked against the 873 training traces by
identifier match, exact normalised text match, and numeric fingerprint combined
with Jaccard overlap. One item matched and was dropped. Neither training set
shares an identifier or problem text with the evaluation set, and the two
training sets are disjoint from each other because they are drawn from
non-overlapping difficulty bands. One limitation is worth stating plainly:
Ganit carries no upstream English source identifiers, so fingerprint and
Jaccard matching are proxies, and a sufficiently free paraphrase could survive
both filters.


\subsection{Evaluation Set}

The held-out set is 500 problems sampled from GANIT-DEV with seed 42,
stratified by difficulty: 145 easy, 127 medium, 129 hard and 99 olympiad. Gold
answers are short final values rather than worked solutions.

Compute constraints on a single consumer GPU limited scoring to the first 250
items of that shuffled set for every model, and to the first 100 for TituLLM-1B
base, which was withdrawn from the extension after producing degenerate output
at 0 to 2\% accuracy across its first hundred items. Every number in this paper
is therefore reported on 250 items unless stated otherwise. Because the set was
shuffled before truncation, the scored subset preserves the difficulty mix: 69
easy, 69 medium, 67 hard and 45 olympiad. All models see the same items in the
same order, so every before-and-after comparison in this paper is paired.

\subsection{Fine-Tuning Protocol}

All models were fine-tuned with LoRA \parencite{hu2021lora} using an identical
configuration: rank 16, $\alpha=32$, dropout 0.05, applied to the attention and
MLP projection matrices (\texttt{q\_proj}, \texttt{k\_proj}, \texttt{v\_proj},
\texttt{o\_proj}, \texttt{gate\_proj}, \texttt{up\_proj}, \texttt{down\_proj}).
Training used a learning rate of $2\times10^{-4}$ with cosine decay and 3\%
warmup, three epochs, per-device batch size 1 with gradient accumulation 16 for
an effective batch of 16, bfloat16 precision, gradient checkpointing, and seed
42. Loss was computed on the target tokens only, with prompt tokens masked.

The training target is the full Bengali reasoning trace from the Ganit
\texttt{messages} field, not the short final answer. Training on the bare answer
would teach the model to emit a value and stop, which is close to the opposite
of reasoning elicitation.

The Bengali instruction text is byte-identical for every model. Only the wrapper
around it differs: completion-style models receive the instruction followed by a Bengali
solution cue (\emph{somadhan}, meaning ``solution''), while instruction-tuned
models receive it through their own chat template. This choice was made empirically rather than assumed. We probed
both formats on each model and selected whichever produced a lower token-cap
rate, on the reasoning that a format a model was never trained in prevents it
from terminating. For Qwen3 models, thinking mode was disabled so that every
backbone receives the same intervention.

All experiments ran on a single NVIDIA RTX 5060 Ti with 16GB of VRAM. Total
training time across all adapters was approximately 4.4 GPU-hours.

\subsection{Availability and Computational Environment}

All experiments ran under Python 3.11.16 with PyTorch 2.8.0+cu128,
Transformers 4.57.6 and PEFT 0.20.0, on a single NVIDIA RTX 5060 Ti with 16GB
of VRAM. Seed 42 was set for Python, NumPy and PyTorch in every run, and all
generation is greedy, so no sampling randomness enters the results.
Deterministic CUDA algorithms were not enabled, because the throughput cost was
prohibitive on this hardware; the consequence is that training losses reproduce
to roughly three decimal places rather than exactly, while greedy generations
are stable.

The curation notebook, the training and evaluation harness, the shared answer
scorer with its executable test suite, and the analysis scripts that regenerate
every table, figure and reported number in this paper from the stored model
outputs are available at [REPOSITORY URL]. The exact environment record,
including package versions and the reproducibility caveats above, is committed
as \texttt{environment.json}.

\subsection{Evaluation Protocol}

Generation is greedy, with a repetition penalty of 1.1 and a per-model
generation budget set to the 99th percentile of that model's own training-target
length. This last point matters and is not a detail. An equal token budget is
not an equal opportunity when tokenizers differ by a factor of three. The same
873 Bengali traces require a median of 397 tokens under TituLLM's tokenizer, 428
under TigerLLM's, and 1,245 under Qwen3's. A uniform 1,024-token cap would leave
the Bengali models room to spare while truncating 42\% of what Qwen3 was trained
to produce. Qwen3 models therefore receive 2,048 tokens and the Bengali models
1,024, with the budget recorded for every prediction.

Generation halts on any end-of-sequence marker the tokenizer knows, rather than
only its declared \texttt{eos\_token}. This is necessary because TituLLM
declares \texttt{<|eot\_id|>}, a chat turn terminator, while being prompted in
completion format where that token never follows free text. Under a
single-token stopping rule the model has no reachable way to stop.

\subsection{Answer Extraction and Scoring}

Answers are scored by exact match after normalisation. Extraction proceeds by
narrowing to an answer region, using \texttt{<answer>} tags, then
\texttt{\textbackslash boxed\{\}}, then the Bengali answer marker \emph{uttar}
(``answer''), and finally the last numeral in the generation, before taking the
first numeric token inside that region. Normalisation maps Bengali numerals
to ASCII, strips currency, percentage and LaTeX markup, and compares numerically
where possible.

Taking the first number inside the answer region rather than the whole region is
essential. An earlier version returned the entire remainder of the line after
the answer marker. A correct answer written naturally, for example
\emph{uttar}: 20 \emph{jon chhatro} (``answer: 20 students''), was compared in
full against a gold value of 20, failed to normalise, and was scored wrong even
though the model was right. The scorer is distributed with an executable
test suite built from verbatim model outputs rather than synthetic cases.

Generations that produce no parseable answer are recorded as format failures and
reported separately from wrong answers.

\subsection{Statistical Analysis}

Because every model is evaluated on the same items, before-and-after comparisons
are paired and use the exact McNemar test on discordant pairs. Comparisons
between different models use Fisher's exact test. Proportions are reported with
Wilson score intervals, which behave correctly near zero where the normal
approximation does not. Continuous measures such as generation length and
Bengali character ratio are reported as mean paired differences with 95\%
confidence intervals, and treated as significant when the interval excludes
zero.

Every before-and-after comparison is reported with a paired risk difference
$\Delta$ and its 95\% confidence interval alongside the $p$-value, because a
$p$-value alone says whether an effect is distinguishable from zero but not how
large it could plausibly be. Where the two disagree we treat the exact test as
authoritative: the Wald interval on a paired difference is anti-conservative
when the discordant counts are very small, which happens here for
TituLLM-1B-Instruct with eight discordant pairs.

We report significance tests rather than raw differences throughout because
accuracies in this study are low. At 3\% accuracy on 250 items, a difference of
three correct answers looks like a trend and is not one.

No correction for multiple comparisons is applied. Seven McNemar tests, three
Fisher tests and seven sign tests are reported. Since every McNemar result is
null, a correction could only make them more null and would not change any
conclusion; the Fisher and sign-test results survive Bonferroni correction at
$\alpha=0.05$ by several orders of magnitude. Readers should nonetheless treat
the single borderline value, $p=0.070$ for TituLLM-1B-Instruct, as weaker than
it looks in isolation.


% =====================================================================
\section{Results}

\subsection{RQ1: Scale versus Targeted Depth}

Table~\ref{tab:main} reports accuracy for every backbone in both conditions.

\begin{table}[htbp]
\caption{Held-out accuracy before and after LoRA elicitation}
\label{tab:main}
\centering
\begin{tabular}{|l|c|c|c|c|}
\hline
\textbf{Model} & \textbf{Bengali PT} & \textbf{Params} & \textbf{Before} & \textbf{After} \\
\hline
TituLLM-1B-Instruct & 37B & 1.2B & \NUM{acc} & \NUM{acc} \\
TituLLM-3B-Instruct & 37B & 3.2B & \NUM{acc} & \NUM{acc} \\
TigerLLM-1B         & 10M & 1.0B & \NUM{acc} & \NUM{acc} \\
Qwen3-0.6B          & --  & 0.6B & \NUM{acc} & \NUM{acc} \\
Qwen3-1.7B          & --  & 1.7B & \NUM{acc} & \NUM{acc} \\
Qwen3-4B            & --  & 4.0B & \NUM{acc} & \NUM{acc} \\
\hline
\end{tabular}
\end{table}

\TODO{Final n=250 numbers, Wilson intervals, and Fisher p-values for the scale
ladder and the depth contrast.}

Two patterns hold. Reading down the columns, accuracy tracks parameter scale.
The Qwen3 ladder rises monotonically from 0.6B to 4B, and each step is
significant in both conditions \NUM{p-values}. Reading across the depth
comparison, a $3{,}700\times$ difference in targeted Bengali pretraining between
TigerLLM and TituLLM produces no significant difference \NUM{p}. Qwen3-4B, which
received no Bengali-specific pretraining at all, outperforms the strongest
Bengali-adapted model by a factor of \NUM{ratio} \NUM{p}.

This answers RQ1 directly, and against our hypothesis. We predicted that
targeted depth would matter more than scale. The data says the opposite.

\subsection{RQ1a: Elicitation Produced No Accuracy Gain}

Reading across the rows of Table~\ref{tab:main}, no model shows a significant
accuracy change from fine-tuning. \TODO{Per-model McNemar p-values.}

This includes Qwen3-4B, which is the important case. LIMO's precondition is that
the base model already holds dense, relevant knowledge that fine-tuning surfaces.
At \NUM{acc} baseline accuracy that precondition is clearly satisfied for
Qwen3-4B, and elicitation still produced no gain.

\subsection{RQ1b: Elicitation Changed Form, Not Capability}

While accuracy did not move, generation behaviour changed substantially and
significantly. Table~\ref{tab:behaviour} reports the behavioural measures.

\begin{table}[htbp]
\caption{Behavioural change after elicitation. Significance from paired 95\% CIs.}
\label{tab:behaviour}
\centering
\begin{tabular}{|l|c|c|c|}
\hline
\textbf{Model} & \textbf{Cap rate} & \textbf{Mean tokens} & \textbf{Bengali ratio} \\
\hline
TigerLLM-1B         & \NUM{b/a} & \NUM{b/a} & \NUM{b/a} \\
TituLLM-1B          & \NUM{b/a} & \NUM{b/a} & \NUM{b/a} \\
TituLLM-1B-Instruct & \NUM{b/a} & \NUM{b/a} & \NUM{b/a} \\
TituLLM-3B-Instruct & \NUM{b/a} & \NUM{b/a} & \NUM{b/a} \\
Qwen3-0.6B          & \NUM{b/a} & \NUM{b/a} & \NUM{b/a} \\
Qwen3-1.7B          & \NUM{b/a} & \NUM{b/a} & \NUM{b/a} \\
Qwen3-4B            & \NUM{b/a} & \NUM{b/a} & \NUM{b/a} \\
\hline
\end{tabular}
\end{table}

Three effects stand out.

First, elicitation taught models to stop. TituLLM-1B never emitted an
end-of-sequence token in any evaluated item before fine-tuning, producing
degenerate repetition until truncation on 100\% of generations. After
fine-tuning it terminated on \NUM{pct} of items and its format failure rate fell
from \NUM{pct} to \NUM{pct}. It solved no additional problems.

Second, the direction of the change depends on scale. \TODO{Qwen3-0.6B
lengthened while Qwen3-1.7B shortened; confirm at final n and give the paired
CIs.} Identical training data moved two models of the same family in opposite
directions.

Third, and least expected, Bengali character adherence fell. \NUM{k} of seven
models showed a significant decline, the largest being \NUM{value}. An
intervention built from Bengali reasoning traces, intended to elicit Bengali
reasoning, made models measurably less Bengali while making them no more
correct. We return to this in Section~\ref{sec:discussion}.

\subsection{Testing the Difficulty-Calibration Explanation}
\label{sec:difficulty}

The most obvious objection to a null result of this kind is that the training
data was simply too hard. Our LIMO-window training set is 74.6\% olympiad
difficulty, and that window was defined by a 32B model's pass rate, not by the
ability of the 0.6B to 4B models actually being trained.

We tested this directly. First we measured whether the models fit their own
training data at all, by evaluating Qwen3-1.7B on 50 problems drawn from its own
training set. The olympiad-trained adapter scored \NUM{k}/50 against \NUM{k}/50
for the base model, so it had barely learned the material it was trained on.

We then retrained the same backbone, with an identical configuration, on the
easy training set described in Section~3.2, and evaluated both adapters on the
same held-out items. Training-set accuracy rose from \NUM{k}/50 to \NUM{k}/50,
confirming that the model can fit properly calibrated data. Held-out accuracy,
however, fell. Table~\ref{tab:difficulty} reports the three-way comparison.

\begin{table}[htbp]
\caption{Effect of training-data difficulty, Qwen3-1.7B, shared held-out items}
\label{tab:difficulty}
\centering
\begin{tabular}{|l|c|c|}
\hline
\textbf{Condition} & \textbf{Accuracy} & \textbf{95\% CI} \\
\hline
Base model, no LoRA      & \NUM{acc} & \NUM{ci} \\
Olympiad-trained adapter & \NUM{acc} & \NUM{ci} \\
Easy-trained adapter     & \NUM{acc} & \NUM{ci} \\
\hline
\end{tabular}
\end{table}

\TODO{Final McNemar p-values for the three paired contrasts, and the
per-difficulty breakdown showing where the easy-trained adapter loses most.}

Successful in-distribution learning transferred negatively. The difficulty
explanation is therefore not available: the models can learn calibrated data,
and doing so does not improve, and may harm, held-out reasoning.

\subsection{The Bengali Tokenizer Tax}

Table~\ref{tab:tokenizer} reports token counts for the same 873 training traces
under each tokenizer.

\begin{table}[htbp]
\caption{Tokens required for identical Bengali reasoning traces}
\label{tab:tokenizer}
\centering
\begin{tabular}{|l|c|c|c|c|}
\hline
\textbf{Tokenizer} & \textbf{Vocab} & \textbf{Median} & \textbf{p95} & \textbf{Max} \\
\hline
TituLLM-1B  & 170,497 & 397   & 873   & 1,500 \\
TigerLLM-1B & 262,145 & 428   & 933   & 1,434 \\
Qwen3       & 151,669 & 1,245 & 2,396 & 3,104 \\
\hline
\end{tabular}
\end{table}

Qwen3 requires roughly $3.1\times$ more tokens than TituLLM for identical
Bengali text, close to the $3.5\times$ reported by
\textcite{yong2025crosslingual} and consistent with the per-model premiums in
\textcite{petrov2023language}. % NEW

The consequence is practical rather than abstract. On a 16GB card, this cost
determined which configurations were trainable at all. Qwen3-4B could not be
trained at a 4,096-token sequence limit and required a reduction to 2,048, which
truncated 104 of 873 traces. No Bengali-tokenized model faced that constraint.
It is worth noting that TigerLLM's efficiency is inherited rather than earned:
Gemma3 ships a 262k vocabulary described as balanced for non-English languages
\parencite{gemma3}, % NEW
and TigerLLM adds exactly one token to it. TituLLM, by contrast, extended
Llama-3.2's vocabulary deliberately for Bengali.


% =====================================================================
\section{Discussion}
\label{sec:discussion}

\subsection{What the Scale Result Means}

Our headline result contradicts TigerLLM's central claim, using TigerLLM's own
released model on a Bengali benchmark. That claim was that targeted,
high-quality pretraining can substitute for scale. We find that it does not, at
least for mathematical reasoning: a $3{,}700\times$ difference in Bengali
pretraining produced no measurable difference, while parameter scale produced a
large and monotonic one.

This is consistent with a literature that has been converging on the same point
from other directions. \textcite{kaplan2020scaling} found performance depends % NEW
strongly on scale and weakly on architectural choices. \textcite{wei2022emergent} % NEW
place the emergence of three-digit arithmetic at roughly 13B parameters and the
benefit of chain-of-thought on GSM8K at roughly 100B. MathMist
\parencite{sobhani2026mathmist} reports the Bengali-English gap closing as Qwen
scales. Our contribution is to separate the variables that this prior work left
entangled, in the one language where the substitution claim was made.

We are careful about what this does not show. It does not show that targeted
Bengali pretraining is worthless. TituLLM's extended tokenizer delivers a real,
measured threefold efficiency gain, and that gain determines what can be trained
on accessible hardware. It shows that targeted depth did not buy mathematical
reasoning.

\subsection{Why Elicitation Changed Form and Not Capability}

LIMA's Superficial Alignment Hypothesis \parencite{zhou2023lima} % NEW
holds that a model's knowledge comes almost entirely from pretraining and that
alignment mainly teaches formatting and style. Our results are a clean, if
unflattering, confirmation of exactly that. Fine-tuning taught these models to
open a reasoning block, to structure a worked solution, and in TituLLM's case to
stop generating at all. It did not teach them to arrive at correct answers.

Two mechanisms plausibly contribute, and we cannot fully separate them.

The first is capacity of the adaptation method. \textcite{biderman2024lora} % NEW
report that LoRA at rank 16, the configuration used here, reaches 0.158 on
GSM8K after continued pretraining where full fine-tuning reaches 0.293, and
conclude that low-rank settings substantially underperform full fine-tuning on
tasks requiring new capability. The first half of that description matches what
we observe.

The second half does not, and the difference is worth stating. Biderman et al.
also report that LoRA forgets less than full fine-tuning and stays closer to
base behaviour. Our retention data in Section~\ref{sec:retention} shows nothing
of the kind: no backbone kept even half of the items it had answered correctly,
and Qwen3-4B abandoned its own output format on 222 of 250 items. We did not
run full fine-tuning, so we cannot claim LoRA forgets more than the alternative.
What we can say is that rank 16 on 873 examples was more than sufficient to
overwrite a model's answering behaviour, and that the common intuition of a
low-rank adapter as a small, conservative edit does not describe this setting.

The second is the capability of the backbones themselves. If arithmetic
competence emerges above 13B \parencite{wei2022emergent}, % NEW
then LIMO's precondition may simply not hold in the 0.6B to 4B range regardless
of adaptation method. Qwen3-4B complicates this reading, since its baseline
accuracy of \NUM{acc} shows real latent capability, and elicitation still
produced nothing.

\subsection{The Bengali Adherence Decline}

The finding we did not anticipate is that training on Bengali reasoning traces
made models less Bengali. \TODO{Confirm count and magnitudes at final n.}

The likely cause is in the traces themselves. Ganit's reasoning contains LaTeX,
ASCII numerals and English mathematical notation, and models trained on it learn
to write mathematics the way the traces do. This connects directly to the
elicitation tax described by \textcite{qi2025xreasoning}, who found that forcing
native Bengali reasoning costs accuracy. Our result is the same tension seen
from the other side: an intervention that did not improve accuracy still moved
models away from Bengali.

For work that aims at genuinely Bengali-native reasoning, this suggests language
adherence should be measured and reported alongside accuracy rather than
assumed. A model that improves on a Bengali benchmark by reasoning in English
has not done what the benchmark intends to measure.


% =====================================================================
\section{Limitations}

\textbf{Adaptation method.} We used LoRA at rank 16 rather than full
fine-tuning. LIMO and s1 both fine-tuned all parameters.
\textcite{biderman2024lora} document that this configuration under-acquires new % NEW
capability, so our null result cannot be cleanly separated from a method
limitation. We report it as a boundary on minimal-data elicitation \emph{under
the adaptation method available at consumer scale}, which is the setting most
researchers in this position face.

\textbf{Training duration.} Models were trained for three epochs. LIMO used
approximately 15 and s1 used 5. Training loss was still decreasing at epoch 3
for four of six backbones, so undertraining cannot be ruled out. No epoch
ablation was run.

\textbf{Architecture is confounded with pretraining depth across families.} The
Bengali-adapted models are built on Gemma3 and Llama-3.2 while the scale ladder
is Qwen3. Comparisons across families are directional rather than causal, as
pre-registered in our study design.

\textbf{One per-model deviation.} Qwen3-4B was trained with a 2,048-token
sequence limit rather than 4,096 because a 4B model at 4,096 exhausted the
available VRAM. This truncated 104 of 873 traces, or 12\%. Every other model
truncated none.

\textbf{Sample size is not uniform.} \TODO{State final per-model n and note
where intervals differ.}

\textbf{The easy subset shifted source composition.} Moving to the easy
difficulty band raised SOMADHAN's share from 13.6\% to 31.0\%, because
GSM8K-derived problems concentrate there. Difficulty and source are therefore
not perfectly separable in that experiment.

\textbf{Training targets retain the source dataset's markup.} Ganit's traces use
\texttt{<think>} and \texttt{<answer>} tags inherited from its Qwen3 lineage. For
Llama and Gemma these are ordinary text tokens, but for Qwen3 \texttt{<think>}
is a control token, and combining disabled thinking mode with a target that
reopens a thinking block is non-canonical. This was applied identically at
training and evaluation, so the comparison is internally valid, but Qwen3 results
may understate properly configured performance.

\textbf{Determinism.} Deterministic CUDA algorithms were not enabled, so losses
reproduce to approximately three decimal places rather than exactly. Reported
differences are orders of magnitude larger than this variance.

\textbf{Scope.} This paper addresses RQ1 only. Cross-domain transfer and
confidence-based stopping were scoped out, as described below.


% =====================================================================
\section{Conclusion}

We set out to determine whether Bengali mathematical reasoning is driven by
parameter scale or by targeted Bengali pretraining, and whether minimal-data
elicitation transfers to model sizes and languages outside the settings where it
was demonstrated.

Scale predicts performance. \TODO{One sentence with the final scale-ladder
figures and p-values.} Targeted Bengali pretraining does not: a
$3{,}700\times$ difference produced no significant effect. This contradicts the
substitution claim made for TigerLLM, tested on that model itself.

Minimal-data elicitation produced no significant accuracy gain on any of seven
backbones spanning 0.6B to 4B parameters, two pretraining philosophies and two
training-difficulty regimes, including a model whose baseline accuracy shows
LIMO's precondition is satisfied. What it reliably produced instead was a change
in form: termination, length, formatting, and a significant decline in Bengali
character adherence.

For researchers working in low-resource languages on accessible hardware, the
practical implication is that the less-is-more result should not be assumed to
carry down in scale. In the range most of us can train, the thousand curated
examples bought presentation, not competence.


% =====================================================================
\section{Future Work}

Three directions follow directly.

\textbf{Adaptation method.} The clearest next experiment is full fine-tuning at
the same scales, or LoRA at substantially higher rank, to separate the method
limitation from the scale boundary. \textcite{biderman2024lora} suggests rank % NEW
256 recovers much of the gap.

\textbf{Cross-domain transfer.} Our original design included a second question:
whether elicited mathematical reasoning transfers to non-mathematical Bengali
tasks. This required the B-REASO benchmark, which was not accessible during the
study period. The question remains open, and TituLLM's own capability-specific
findings \parencite{nahin2025titullm} suggest it is worth asking.

\textbf{Adaptive test-time compute.} A third question we scoped out is whether a
model can judge when it has reasoned enough, using cheap untrained confidence
signals rather than the calibration pipelines that current methods require. The
test-time scaling literature is unsettled on whether longer reasoning helps at
all, and none of it has been tested in Bengali. With a measured $3.1\times$
tokenizer penalty, the cost of unnecessary reasoning tokens is higher in Bengali
than the English results suggest, which makes the question more pressing rather
than less.
